%% Glossary & Abbreviations
\mysection{Glossary \& Abbreviations}

\subsection*{Abbreviations}

\begin{tabular}{ll}
AI          & Artificial Intelligence \\
API         & Application Programming Interface \\
LitM        & Lost in the Middle \\
LLM         & Large Language Model \\
NLP         & Natural Language Processing \\
PLM         & Pre-trained Language Model \\
RAG         & Retrieval-Augmented Generation \\
TTFT        & Time-to-First-Token \\
\end{tabular}

\vspace{0.6cm}

\subsection*{Glossary}

\begin{tabular}{lp{9.5cm}}
\textbf{Chunking}           & The process of dividing a large source text into smaller, distinct segments (chunks) prior to vector database insertion. \\
\textbf{Context Window}     & The maximum amount of text (measured in tokens) that a language model can process in a single prompt iteration. \\
\textbf{Cosine Similarity}   & A mathematical metric used to measure the semantic similarity between two vectors by calculating the cosine of the angle between them. \\
\textbf{Embedding}          & A mathematical vector representation of text that captures its semantic and contextual meaning in a high-dimensional space. \\
\textbf{Hallucination}      & A phenomenon where a language model generates output that is fluent and plausible but factually incorrect or unsupported by the training data. \\
\textbf{Inference}          & The operational phase where a fully trained language model processes a user prompt and generates an output response. \\
\textbf{Non-parametric Knowledge} & Data or information that resides outside the language model's internal weights, typically stored in an external database, and retrieved dynamically during inference. \\
\textbf{Parametric Knowledge} & The static, internal information that is permanently embedded within a language model's weights during its training phase. \\
\textbf{Time-to-First-Token} & A performance metric that measures the latency from the moment a prompt is sent to the model until it begins generating its very first output token. \\
\textbf{Token}              & The foundational atomic unit of text (such as a word, character, or sub-word) into which input strings are split before being processed by a language model. \\
\textbf{Vectorization}      & The algorithmic process of converting textual data into numerical vectors (embeddings) using a machine learning model. \\
\end{tabular}